\input{preamble/preamble.tex}

\geometry{margin=0.85in}
\AtBeginDocument{\small}

\newcommand{\ExamNameField}{\noindent\textbf{Name:}\ \rule{0.7\linewidth}{0.4pt}\par\medskip}

\begin{document}
	\ExamTitleBlock{11th grade}{Learning evidence 2.1: Confidence Interval Planning}{}

	\ExamNameField

	\ExamSection{Evaluated criteria}
	This exam evaluates only the criteria C1--C5. Each criterion is evaluated across the three problems below.
	A criterion is passed if it is correctly activated in at least two of the three problems.
	The overall exam result is binary (Passed/Not passed).

	\begin{ExamCriteria}
		\item C1 Compute the sample mean and sample standard deviation.
		\item C2 Distinguish between sample statistics and population parameters.
		\item C3 Explain why interval estimation is preferred over a single point estimate.
		\item C5 Construct a 95\% confidence interval using known population standard deviation.
		\item C4 Interpret the meaning of a 95\% confidence interval in context.
	\end{ExamCriteria}

	\begin{ExamProblems}
		\item A city budgeting office wants to estimate the average weekly transit subsidy per rider for two pilot zones.
		Zone A recorded the following sample of weekly subsidy amounts (USD):

		\begin{itemize}
			\item Zone A sample (\(n_A = 6\)): 18, 22, 19, 21, 20, 17.
			\item Zone B sample (\(n_B = 8\)): 24, 26, 23, 25, 27, 22, 24, 26.
		\end{itemize}

		An annual audit provides the population standard deviations for the zones: \(\sigma_A = 3.5\) and \(\sigma_B = 4.0\).
		For classroom purposes, suppose the true population mean weekly subsidy is \(\mu_A = 20\) dollars in Zone A and \(\mu_B = 25\) dollars in Zone B.
		Construct and interpret a 95\% confidence interval for the population mean weekly subsidy in each zone.

		\item A regional warehouse studies the number of orders processed per hour to plan staffing.
		A random sample of 60 hours is summarized in grouped form below. The operations team also reports a known population standard deviation of \(\sigma = 6.2\) orders per hour.
		For academic purposes only, assume the true population mean order rate is \(\mu = 22\) orders per hour.
		The grouped data pairs (orders per hour, frequency) are: \(14 \rightarrow 8\), \(18 \rightarrow 16\), \(22 \rightarrow 20\), \(26 \rightarrow 10\), and \(30 \rightarrow 6\).
		Construct and interpret a 95\% confidence interval for the population mean orders per hour.

		\item A logistics firm monitors the time (in minutes) required to load containers at a port.
		A sample of 80 loading times is grouped into class intervals below. The firm has a historical estimate of the population standard deviation, \(\sigma = 5.8\) minutes.
		For classroom reference, suppose the true population mean loading time is \(\mu = 32\) minutes.
		The grouped intervals and frequencies are: 20--24 minutes (12), 25--29 minutes (18), 30--34 minutes (20), 35--39 minutes (17), and 40--44 minutes (13).
		Construct and interpret a 95\% confidence interval for the population mean loading time.

	\end{ExamProblems}

\ifdefined\ExamBook
\else
\end{document}
\fi
